\chapter{Quarter-Disk Zero Exclusion and the Complete Negative-Inversion Spectral No-Go}
\label{chap:quarter-disk-zero-exclusion}

\status{SOH-G018 proved: the closed quotient disk $|w|\le1/4$ is zero-free; consequently the canonical negative inversion maps no zero of $F$ or $\xi$ to another zero.}

\section{Purpose}

The preceding sequence G014--G017 progressively constrained the possible spectral action of the canonical negative inversion. G014 and G015 proved that paired zero sets are finite, G016 identified their exact two-to-one quotient correspondence, and G017 excluded all fixed exceptional roots. This chapter closes the remaining finite two-cycle possibility.

The argument uses only the positive Taylor coefficients already proved for the even quotient
\[
 \xi\!\left(\frac12+z\right)=F(z^2),
 \qquad F(w)=\sum_{k\ge0}a_k w^k,
 \qquad a_k>0,
\]
together with an elementary rational lower bound for $F(0)=\xi(1/2)$.

\section{An elementary lower bound for $\xi(1/2)$}

With the standard completed-xi normalization,
\[
 \xi\!\left(\frac12\right)
 =\frac18\pi^{-1/4}\Gamma\!\left(\frac14\right)
 \bigl(-\zeta(1/2)\bigr).
\]

The alternating Dirichlet eta series yields
\[
 \eta\!\left(\frac12\right)
 >1-\frac1{\sqrt2}+\frac1{\sqrt3}-\frac12.
\]
The exact integer inequalities
\[
 2\cdot71^2>100^2,
 \qquad
 3\cdot57^2<100^2
\]
give
\[
 \frac1{\sqrt2}<\frac{71}{100},
 \qquad
 \frac1{\sqrt3}>\frac{57}{100},
\]
and therefore
\[
 \eta\!\left(\frac12\right)>\frac9{25}.
\]
Since
\[
 \eta(1/2)=(1-\sqrt2)\zeta(1/2),
\]
we have
\[
 -\zeta(1/2)=(\sqrt2+1)\eta(1/2).
\]
Using $\sqrt2>7/5$ gives
\[
 -\zeta(1/2)>\frac{108}{125}.
\]

For the Gamma factor, on $0\le t\le1$ the alternating Taylor remainder gives
\[
 e^{-t}\ge1-t+\frac{t^2}{2}-\frac{t^3}{6}.
\]
Hence
\[
 \Gamma\!\left(\frac14\right)
 >\int_0^1 t^{-3/4}
 \left(1-t+\frac{t^2}{2}-\frac{t^3}{6}\right)dt
 =\frac{1972}{585}.
\]

Finally, the classical bound $\pi<22/7$ and the exact comparison
\[
 \frac{22}{7}<\left(\frac{50}{37}\right)^4
\]
imply
\[
 \pi^{-1/4}>\frac{37}{50}.
\]
Combining the three bounds gives
\[
 \boxed{
 \xi\!\left(\frac12\right)
 >\frac18\frac{37}{50}\frac{1972}{585}\frac{108}{125}
 =\frac{54723}{203125}
 >\frac14.}
\]

\section{The quarter-disk theorem}

At the positive boundary point,
\[
 F\!\left(\frac14\right)=\xi(1)=\frac12.
\]
Because all Taylor coefficients are positive, for every $|w|\le1/4$,
\[
 \begin{aligned}
 |F(w)|
 &\ge a_0-\sum_{k\ge1}a_k|w|^k\\
 &\ge a_0-\sum_{k\ge1}a_k4^{-k}\\
 &=2a_0-F(1/4).
 \end{aligned}
\]
Using the lower bound for $a_0=F(0)=\xi(1/2)$ gives the uniform estimate
\[
 \boxed{
 |F(w)|>
 2\frac{54723}{203125}-\frac12
 =\frac{15767}{406250}>0
 \qquad(|w|\le1/4).}
\]
Therefore the complete closed disk is zero-free:
\[
 \boxed{F(w)\ne0\quad\text{whenever }|w|\le1/4.}
\]
Every zero of $F$ consequently satisfies $|w|>1/4$.

\section{Complete spectral no-go for negative inversion}

The quotient negative inversion is
\[
 J(w)=\frac1{16w}.
\]
If $F(w)=0$, then $|w|>1/4$, so
\[
 |J(w)|=\frac1{16|w|}<\frac14.
\]
The quarter-disk theorem then gives
\[
 F(J(w))\ne0.
\]
Thus the paired quotient set is empty:
\[
 \boxed{P_J=\varnothing},
 \qquad
 \boxed{Z_F\cap J(Z_F)=\varnothing}.
\]

G016 proved that the quotient map
\[
 q(s)=\left(s-\frac12\right)^2
\]
identifies the paired xi-zero set with $P_J$. Hence
\[
 \boxed{P_N=\varnothing},
 \qquad
 \boxed{Z_\xi\cap N_s(Z_\xi)=\varnothing}.
\]
The canonical negative inversion therefore remains an exact projective and spinorial operator of the coordinate geometry while having no zero-to-zero spectral action on the Riemann xi function.

\section{What this closes}

The chain is now exact:
\begin{itemize}
 \item G014: the paired xi-zero set is finite;
 \item G015: the paired quotient-root set is finite;
 \item G016: the two sets are related exactly two-to-one;
 \item G017: neither fixed point of the quotient involution is a root;
 \item G018: both paired sets are empty.
\end{itemize}

No table of Riemann zeros enters the proof. Numerical boundary sampling retained in the software is only a regression test against the analytic margin.

\section{Proof firewall}

G018 proves the rational lower bound for $\xi(1/2)$, the uniform quarter-disk lower bound for $|F|$, the zero-free disk $|w|\le1/4$, and the complete absence of negative-inversion paired zeros. It does not prove real-rootedness of $F$, PF$\infty$, SOH-G003 real-rootedness, or the Riemann Hypothesis.
